\documentclass[12pt,letterpaper]{article}

%Packages
\usepackage{pgfplots}
\usepackage{caption}
\usepackage{subcaption}

\usepgfplotslibrary{external}
\pgfplotsset{width=5cm,compat=1}
\usepackage{hyperref}
\usepackage{amsmath}
\usepackage{amsthm}
\usepackage{amsfonts}
\usepackage{amssymb}
\usepackage{amscd}
\usepackage{dsfont}
\usepackage{enumerate}
\usepackage{fancyhdr}
\usepackage{mathrsfs}
\usepackage{bbm}
\usepackage{framed}
\usepackage{mdframed}
\usepackage{cancel}
\usepackage{float}
\usepackage{mathtools}
%\usepackage[]{mcode}
\usepackage{graphicx}

%Page formatting
\usepackage[letterpaper,voffset=-.5in,bmargin=3cm,footskip=1cm]{geometry}
\setlength{\parindent}{0.0in}
\setlength{\parskip}{0.1in}
\allowdisplaybreaks
\headheight 15pt
\headsep 10pt

% Common Commands
\newcommand\N{\mathbb N}
\newcommand\Z{\mathbb Z}
\newcommand\R{\mathbb R}
\newcommand\Q{\mathbb Q}
\newcommand\lcm{\operatorname{lcm}}
\newcommand\setbuilder[2]{\ensuremath{\left\{#1\;\middle|\;#2\right\}}}
\newcommand\E{\operatorname{E}}
\newcommand\V{\operatorname{V}}
\newcommand\Pow{\ensuremath{\operatorname{\mathcal{P}}}}

\DeclarePairedDelimiter\ceil{\lceil}{\rceil}
\DeclarePairedDelimiter\floor{\lfloor}{\rfloor}

\newcommand\definition{\textbf{Definition: }}
\newcommand\proposition{\textbf{Proposition: }}

% 22 style
\newcommand\hint[1]{\textbf{Hint}: #1}
\newcommand\note[1]{\textbf{Note}: #1}
\newenvironment{22enumerate}{\begin{enumerate}[a.]\itemsep0em}{\end{enumerate}}
\newenvironment{22itemize}{\begin{itemize}\itemsep0em}{\end{itemize}}
\fancypagestyle{firstpagestyle} {
  \renewcommand{\headrulewidth}{0pt}%
  \lhead{\textbf{CSCI 0220}}%
  \chead{\textbf{Discrete Structures and Probability}}%
  \rhead{Littman}%
}

\pagestyle{fancyplain}

\newcommand\EX{\mathds{E}}
\newcommand\PR{\mathds{P}}
\newcommand\zlam{Z_\lambda}
\DeclareMathOperator*{\argmin}{arg\,min}

%%%%%%%%%%%%%%%% COMPILE WITH \soltrue or \solfalse %%%%%%%%%%%%%%%%%%%%%%%%%%%%%
\newif\ifsol
\soltrue  % SOLUTIONS
% \solfalse % NO SOLUTIONS
%%%%%%%%%%%%%%%%%%%%%%%%%%%%%%%%%%%%%%%%%%%%%%%%%%%%%%%%%%%%%%

%%%%%%%%%%%%%%%% solution help %%%%%%%%%%%%%%%%%%%%%
\newcommand{\solu}[2]{ \begin{mdframed} \ifsol #2 \else \vspace{#1} \fi \end{mdframed} }
\newcommand{\solt}{ \ifsol (T) \fi }
\newcommand{\solf}{ \ifsol (F) \fi }
\newcommand{\solm}[1]{\ifsol \textit{(#1)} \fi}


\begin{document}
  \thispagestyle{firstpagestyle}
  \begin{center}
    {\large \textbf{Recitation 4}}
    
    {\large Relations and Functions}

  \end{center}

\section*{Part 1: Relations}

	\subsection*{Definitions}

\textbf{Defn 1:} A \textit{relation} $R:A \rightarrow B$ is defined on a domain $A$, co-domain $B$, and a graph that is a subset of the Cartesian product $A \times B$. As a slight abuse of notation, we will often write the relation $R$ as a subset of $A\times B$. A relation \empn{on} a set $A$ is $R:A \rightarrow A$. 

\textbf{Defn 2:} An \textit{equivalence relation} is a relation that is reflexive, symmetric, and transitive.

\textbf{Defn 4:} A relation $R$ on $A$ is \texit{reflexive} if $\forall a \in A$, $(a, a) \in R$.

\textbf{Defn 5:} A relation $R$ on $A$ is \textit{symmetric} if $\forall (a, b) \in R$, $(b, a) \in R$. An equivalent definition is that a relation is \textit{not} symmetric if $ \exists (a, b) \in R$ such that $(b, a) \notin R$.

\textbf{Defn 6:} A relation $R$ on $A$ is \textit{antisymmetric} if $\forall a, b \in A$, $(a, b) \in R$ and $(b, a) \in R$ implies that $a = b$.

\textbf{Defn 7:} A relation $R$ on $A$ is \textit{transitive} if $\forall (a, b), (b, c) \in R$, $(a, c) \in R$. An equivalent definition is that a relation is \textit{not} transitive if $ \exists (a, b), (b, c) \in R$ such that $(a, c) \notin R$.

\textbf{Defn 8:} Let $R$ be an equivalence relation on $A$. Then, the \textit{equivalence class} of $a \in A$, denoted $\left[a\right]_R$,
is $\{x\in A \mid (a, x) \in R\}$. That is, 
$\left[a\right]_R$ is all of the elements to which $a$ is related.

\subsection*{What is the point of an equivalence relation, anyway?}

What does it mean for two things to be equal? It can depend on context. For example, you probably generally think of the numbers 2 and 4 as not being equal. However, maybe I want to consider the numbers 2 and 4 to be equal in some contexts because they are both even. We could be in a situation where we only want there to be two kinds of things: even things and odd things. We don't care about anything else like how big or how small the thing is. 

\textit{An equivalence relation allows us to specify what things in the world are equal to each other, and what things aren't.} % This is why they're called \textbf{equivalence} relations! 

An equivalence relation $R$ splits up 
% (``partitions'') 
our world into categories, or equivalence classes. In a given equivalence class, all things within the class are things we consider equal, or equivalent, in the context of $R$.

For instance, the equivalence relation $R = \{(x, y) \mid x \bmod 2 = y \bmod 2\}$, for all $x, y \in \Z$, divides the world into two categories: the odd and even numbers.

We call the way the equivalence relation splits up our world a \textit{partition}. A little more formally, a \textit{partition} of a set $A$ is a list of subsets $B_1, \ldots, B_k$ of $A$ such that every element of $A$ is in some subset $B_i$ (exhaustive), but no two subsets share an element (mutually exclusive). \\

\includegraphics[scale=.7]{partition.png}  \\
\textit{One possible partition of some set $A$, where the dots in the square are distinct elements of $A$}

\subsection*{Task 1}	

	\begin{enumerate}[1.]
        
		\item Consider the set $A = \{1,2,3\}$. In the following questions, all relations are on $A$. It may be helpful to draw out a diagram of each relation.
		
			\begin{enumerate}[a.]
				\item $R = A \times A$. List out the elements of $R$. Is $R$ an equivalence relation? If so, state its equivalence class(es).

				\solu{2cm}{$\{(1,1), (1,2),(1,3),(2,1),(2,2),(2,3),(3,1),(3,2),(3,3)\}$. It is an equivalence relation because everything is related to each other, since it's the entire Cartesian product. There is only one equivalence class: the entire set.}
				
				
				\item  $R = \{(1, 2), (2, 1)\}$. Is this relation \textit{transitive}?
				
				\solu{1cm}{No, it's not transitive. $(1, 2), (2, 1) \in R \implies (1, 1) \in R$, but it's not actually in $R$. The same argument works for $(2, 2)$.}
				
				\item  $R = \{(1, 2), (2, 1),(2,2),(1,1)\}$. Is this relation reflexive? Symmetric? Transitive?
				
				\solu{3cm}{$(3, 3)$ is not in $R$, so it's not reflexive (since we defined $R$ to be on $A$. However, it is transitive and symmetric.}
				
				\item If the relation in question iii is not an equivalence relation, can you add one pair to it and make it an equivalence relation? Write the equivalence classes of the new relation. 
				
				\solu{2cm}{Yes, add $(3, 3)$. In that case, the equivalence classes become $\{1, 2\}$ and $\{3\}$.}
				
				\end{enumerate}
	\item Let $A=\{1,2\}$ and answer to the following questions.
				\begin{enumerate}[a.]
				
				\item What is the equivalence relation on $A$ with the smallest number of equivalence classes possible?
				
				\solu{1cm}{$A \times A$, a world where everything is equal to everything.}
				
				\item What is the equivalence relation on $A$ with the largest number of equivalence classes possible?
				
				\solu{1cm}{$\{(1, 1), (2, 2)\}$, where things are equal only if they are truly equal.}

				\item Is  $R_0 = \{\}$ a relation on $A$?

				\solu{1cm}{Yes, the empty set is a subset of every set}

				\item Is $R_0$ symmetric? Is it antisymmetric? Why or why not?
				
				\solu{2cm}{Yes, it does not violate our definition of symmetry/antisymmetry as there are no pairs to begin with. (``vacuously true" based on logical equivalence  $(F\implies P)\equiv T$)}

				\item Is $R_0$ transitive? Why or why not?

				\solu{2cm}{Yes, it does not violate our second definition of transitive as there are no pairs. (``vacuously true based on the logical equivalence $(F\implies P)\equiv T$")}

				\item 	$R_0$ is not an equivalence relation. Why?
				
				\solu{1cm}{It is not reflexive.}
		\end{enumerate}

 
    \item Suppose $R$ is an equivalence relation on $S$, and $R = \{\}$. What is $S$? 
        
        \solu{1cm}{The empty set, otherwise it is not reflexive.}
       

		\item Consider the set $B$ of all students at Brown. For each of the following relations on $B$, state whether they are reflexive, symmetric, antisymmetric, transitive, or some combination of them. If it is an equivalence relation, then determine the equivalence classes of the relation.
		
		\begin{enumerate}[a.]
			\item Two students are related if they have the same astrology sign.


				\solu{4cm}{Reflexive, symmetric, and transitive. Therefore equivalence relation. Equivalence classes are students of same astrology sign. 
				
				However, not antisymmetric.}


			\item$s_1$ and $s_2$ are students and $(s_1,s_2) \in R$ if $s_1$ is younger than or the exact same age as $s_2$. (You can assume no students were born at the exact same time.)

				\solu{4cm}{Transitive and anti-symmetric and reflexive, but not symmetric.}

		\item Two students are related if they are studying anthropology.
		
				\solu{4cm}{Symmetric and transitive but not reflexive or anti-symmetric.}


	  	\item Two students are related if they go to Brown.


				\solu{4cm}{Reflexive, symmetric, and transitive, but not antisymmetric. Therefore equivalence relation. One equivalence class which consists of all students at Brown.}
		\end{enumerate}
    
\end{enumerate}

\textbf{Checkpoint 1 — Call a TA over!}

\newpage

	\section*{Part 2: Functions}

      \subsection*{Definitions}

      \textbf{Defn 1:} A relation $R: X\rightarrow Y$ is a \textbf{function} if for every $x$ in the domain $X$, $x$ is mapped to one and only one $y$ in $Y$, the codomain. Note that in the book this is called a \textit{total function}, and function refers to a \textit{partial function}, where for every $x$ in the domain $X$, $x$ is mapped to zero or one $y$ in the codomain $Y$. In this class, we will use function to mean total function and partial function to mean partial function.

      \textbf{Defn 2:} The \textbf{range} of a function $f$ consists of all members of the codomain of $f$ that are mapped to by some member of the domain of $f$. It is the \emph{image} of the domain.

      \textbf{Defn 3:} $f : X \rightarrow Y$ is \textbf{injective (one-to-one)}  if, for every $y \in Y$, there is \textit{at most one} $x \in X$ such that $f(x) = y$. Equivalently, for any $x,y\in Y$ we have $ f(x) = f(y)\implies x=y$, and you can also use its contrapositive $x_1 \not= x_2 \implies f(x_1) \not= f(x_2)$.
    %   We can also define this property on all relations instead of just functions.

      \textbf{Defn 4:} $f : X \rightarrow Y$ is \textbf{surjective (onto)} if, for every $y \in Y$, there is at least one $x \in X$ such that $f(x) = y$. For surjective functions, the range is equal to co-domain.
    %   We can also define this property on all relations instead of just functions.

      \textbf{Defn 5:} $f : X \rightarrow Y$ is a \textbf{bijection} if it is both an injection and surjection. 

      \subsection*{Task 2}
      
      Let $A$ be the set $\{1, 2, 3\}$. Consider the following relation on $A$, $R_1$ = $\{(1, 2), (2, 1)\}$. 

      \begin{enumerate}[1.]
      \item Is $R_1$ a function? 
      \solu{1 cm}{No; not all members of $A$ are mapped to something in $A$. It is a partial function but not a (total) function.}
      \end{enumerate}

      Now, consider $R_2$, another relation on $A$:  $\{(1, 2), (2, 1), (3, 2)\}$. 
      \begin{enumerate}[1.]
      \item Is $R_2$ a function? 
      \solu{1 cm}{Yes.}
      \item If $R_2$ is a function, what is its codomain? How about its range?
       \solu{1 cm}{The codomain is $A$. The range is $\{1, 2\}$.}
      \end{enumerate}

      \subsection*{Task 3}
      
      Consider these diagrams that visualize a relation $R:A \rightarrow B$. The diagrams have two sets of dots, one for $A$ and one for $B$, and they have an arrow from $a$ to $b$ in whenever $(a, b)\in R$.
      
      \includegraphics[scale=.5]{diagrams.png}
      
      Match each of the five diagrams, labeled \sf{A}--\sf{E}, with one of these five descriptions below:
      \begin{enumerate}[1.]
          \item \underline{\hspace{5mm}} Not a function
          \item \underline{\hspace{5mm}} A function that is neither surjective nor injective
          \item \underline{\hspace{5mm}} A surjective function that is not injective
          \item \underline{\hspace{5mm}} An injective function that is not surjective
          \item \underline{\hspace{5mm}} A bijective function — both surjective and injective
      \end{enumerate}
      
      \solu{3 cm}{
        \begin{enumerate}[1.]
          \item E is not a function
          \item A is a function that is neither surjective nor injective
          \item C is a surjective function that is not injective
          \item B is an injective function that is not surjective
          \item D is a bijective function
      \end{enumerate}
      }

\textbf{\textit{Optional} Checkpoint (recommended if queue is short) — Call a TA over!}

\newpage
\subsection*{Task 4}
 

 Consider the following functions and determine if the given function is an injection, surjection, and/or bijection.

\begin{figure}[ht]
 \begin{subfigure}[b]{0.3\textwidth}
\begin{tikzpicture}
  \draw[->] (-2, 0) -- (2, 0) node[right] {$x$};
  \draw[->] (0, -2) -- (0, 2) node[above] {$y$};
  \draw[scale=0.3, domain=-3:3, smooth, variable=\x, blue] plot ({\x}, {\x*\x});
\end{tikzpicture}
 \caption{ $f(x)=x^2$}
\end{subfigure}
\hfill
 \begin{subfigure}[b]{0.3\textwidth}
\begin{tikzpicture}
  \draw[->] (-2, 0) -- (2, 0) node[right] {$x$};
  \draw[->] (0, -2) -- (0, 2) node[above] {$y$};
  \draw[scale=0.3, domain=-3:3, smooth, variable=\x, blue] plot ({\x}, {2*\x});
\end{tikzpicture}   
 \caption{ $g(x)=x/2$}
 \end{subfigure}
 \hfill
  \begin{subfigure}[b]{0.3\textwidth}
\begin{tikzpicture}
  \draw[->] (-2, 0) -- (2, 0) node[right] {$x$};
  \draw[->] (0, -2) -- (0, 2) node[above] {$y$};
  \draw[scale=0.3, domain=-2:2, smooth, variable=\x, blue] plot ({\x}, {\x*(\x*\x-1)});
\end{tikzpicture}   
 \caption{ $h(x)=x^3-x$}
 \end{subfigure}
 \end{figure}
 
      Discuss your answers!
      \begin{enumerate}[a.]
      
		
		\item $f : \R \rightarrow \R$, $f(x) = x^2$

		\solu{1 cm}{Not injective or surjective.}	
		
		 \item $g : \R \rightarrow \R$, $g(x) = \frac{x}{2}$

		\solu{1cm}{Injective and surjective. Therefore bijective.}
		
        \item $h:\R\rightarrow \R$, $h(x)=x^3-x$
        
        \solu{1cm}{Surjective, but not injective}
        
         \item Question: All of the above functions are defined on $\R$. Consider their graphs in the coordinate system. Which of the following implies surjectivity, which implies injectivity, and which implies neither?
 \begin{enumerate}[(1)]
     \item any vertical line intersects the graph \textbf{at most} once
     \item any horizontal line intersects the graph \textbf{at most} once
      \item any vertical line intersects the graph \textbf{at least} once
     \item any horizontal line intersects the graph \textbf{at least}  once
 \end{enumerate}

\solu{2cm}{Surjection: (4).
Injection: (2). Neither: (1), (3)}

        \item Recall the functions defined in parts a-c. Is $f\circ h:\R\rightarrow \R$ surjective and/or injective? (Use \href{https://www.desmos.com/calculator}{\color{blue} a graphing calculator} if you need to.)
        
        \solu{1cm}{$f\circ h(x)= (x^3-x)^2$, not surjective,  not injective}
        
        \item Is $h\circ f: \R\rightarrow \R$ surjective and/or injective?
        
        \solu{1cm}{$h \circ f(x)= x^6 - x^2$, not surjective,  not injective}
        


        \item Let $f : \R \rightarrow \Z$
        % , where
        % $f(x) = \lfloor x \rfloor $.
        % Let $f$
        give as output the greatest integer less than or equal to $x$, denoted as the floor function $f(x) = \lfloor x\rfloor$. For instance, $f(3.5)=3$, $f(3)=3$, and $f(\pi)=3$.\footnote{Note that we can similarly define the ceiling function $f : \R \rightarrow \Z$ that gives as output the smallest integer greater than or equal to $x$, denoted the ceiling function $f(x) = \lceil x \rceil $. For example, $hf(3.5)=4$, $f(3)=3$, and $f(\pi)=4$.}
        
        \solu{1cm}{Surjective, but not injective.}
 
        \item $f: \Z \rightarrow \Z$
        
        $f(x) = \lfloor x \rfloor$

        \solu{1cm}{Injective and surjective. Therefore bijective.}

      %  $f(\text{player}) = \text{the player's role}$

	%	\solu{1cm}{Not injective. Surjective.}		


%         \item \textit{Optional:}
        
%         $f : \{0,1\} \rightarrow \N$

%         $f(0) = 1, f(1) = 0$ 

% 	  \solm{Injective but not surjective.}
        \item \textit{Optional:} $f : \text{Brown University Students} \rightarrow \text{Countries in the World}$

        $f(\text{student}) = \text{country where student is from}$

		\solu{1cm}{Not injective, and sadly not surjective.}
		
        \item \textit{Optional:} $f : \text{First Year Students} \rightarrow \text{First Year Dorms}$ 

        $f(\text{student}) = \text{dorm that student lives in}$

        \solu{1cm}{Not injective. Presumably surjective.}

      \end{enumerate}
      
\subsection*{Task 5}

Let $A=\{0,1,2\}$ and the function $f:{\mathcal P}(A)\rightarrow \{0,1\}^3$.

Denote the output as $(f_0, f_1, f_2)$, where each $f_i$ is 0 or 1. Given an element $X \in \mathcal{P}(A)$, define $f(X)$ where each $f_i$ is 1 iff $i\in X$, and 0 otherwise.
 
 For instance, $f(\{0, 2\}) = (1, 0, 1)$, $f(A) = (1, 1, 1)$, and $f(\emptyset) = (0, 0, 0)$.
 
 \begin{enumerate}[a.]
     \item Consider a set of size $n$ where $A=\{1,2,\ldots, n\}$. Generalize the above bijective function for ${\mathcal P}(A)$ and binary strings of length $n$ (that is, $\{0,1\}^n$).
     
     \solu{2cm}{ $f:{\mathcal P}(A)\rightarrow \{0,1\}^n$, let $X \in \mathcal{P}(A)$, for any $i=1,2, \ldots, n$, the $i$-th element in $f(X)$ is 1 iff $i\in X$.}
     
     \item Prove that the above function is a bijection.
     
     \solu{8cm}{
     \begin{description}
        \item[Surjectivity.] Given any $(f_0, f_1, \ldots, f_n)$, we can construct a subset $X$ of $A$ such that $f(X)$ results in the above binary string. For each element $i \in A$, include it in $X$ iff $f_i = 1$. Then, by the definition of our function, we will result in the exact tuple that we started out with, proving surjectivity.
        
        \item[Injectivity.] We'll show that $X \neq Y \implies f(X) \neq f(Y)$. If the two subsets are not equal, it means that there is at least some element $e$ that one has but the other does not. Without loss of generality, say $e \in X$ but $e \notin Y$. 
        
        As a result, $f_e=1$ for $f(X)$ but $f_e = 0$ for $f(Y)$, so the resulting tuples are not equal. Therefore, we've shown injectivity.
        
        With both properties shown, $f$ is a bijective function.
     \end{description}
     }
 \end{enumerate}

\subsection*{\textit{Optional} Task 6}
      
      \begin{enumerate}[a.]
      \item If a function $f : X \rightarrow Y$ is injective, what can we say about the cardinalities of $X$ and $Y$? Try making some diagrams where $X$ has more elements than $Y$, fewer elements than $Y$, or the same number of elements as $Y$. When are you able to create an injection, and when are you not?
      
      \solu{2 cm}{If and only if $|X| \leq |Y|$, then there exists $f : X \rightarrow Y$ such that $f$ is injective.}
      
      \item If a function $f : X \rightarrow Y$ is surjective, what can we say about the cardinalities of $X$ and $Y$? Again, you might want to draw out some examples.
      
      \solu{2 cm}{If and only if $|X| \geq |Y|$, then there exists $f : X \rightarrow Y$ such that $f$ is surjective.}
      
      \item Based on what you've found in the previous two questions, if a function $f: X \rightarrow Y$ is bijective, what can we say about the cardinalities of $X$ and $Y$? When can we create a bijection between two sets, and when can we not?
      
      \solu{2 cm}{If and only if $|X| = |Y|$, then there exists $f : X \rightarrow Y$ such that $f$ is bijective.}

      \end{enumerate}
      
\textbf{Checkoff - Call a TA over!}

\end{document}